\subsection{Toric cuts with unequal paired lengths}
\label{sec:A1-toric-unequal}

There is a fourth reading of the cut datum which must be distinguished
from the geodesic reading. Require that the identification of the two
banks of every cut be toric for their canonical, normalized arclength
circle actions, but do not require their hyperbolic circumferences to
be equal. Toric compatibility is a condition on normalized angles;
it does not equate these two lengths. In this subsection a cut is
an actual analytic embedded seam, and isomorphisms must preserve the
seams and the external labels. The phase is an actual circle
identification, rather than a fully marked real twist.

Fix a stable graph $\Gamma$, its marked normalization $Y=(Y_v)$,
and a subset $L$ of the external markings which are actual boundary
components; the other external markings remain cusps. Let $H$ be
the set of the $2|E|$ node branches together with $L$, and let
$P=E\sqcup L$. Every branch or leg in $H$ has its own independent
positive squared length parameter
\[
 s_h=(\ell_h/2\pi)^2>0.
\]
At each component choose its unique canonical cap system with these
parameters, as supplied by Theorem~\ref{A4-grafting}. Glue the node
boundaries torically. The analytic collar gluing, equivalently the
analytic case of Theorem~\ref{A3-welding}, produces a Riemann surface
with its actual seams even when the paired lengths differ. Conversely,
cutting such a surface and canonically capping recovers these data.
We denote this space, or the corresponding groupoid when
automorphisms are retained, by $\mathcal R_{\Gamma,L}(Y)$.

Put $S(s)=\sum_v S_v(s|_{J_v})$, where $J_v$ is the active set at
the vertex. Theorem~\ref{A2b-image} identifies the full cap logarithm
map
\[
 F=\nabla S:(0,\infty)^H\longrightarrow D_Y
\]
with a real-analytic diffeomorphism onto an open convex downward
closed subset $D_Y\subset\R^H$. A vertex with no active flags
contributes no variables. In fixed tangent frames define the
surjective linear map
\[
 P_*:\R^H\longrightarrow\R^P,\qquad
 (P_*x)_e=x_{h_e^+}+x_{h_e^-},\quad
 (P_*x)_l=x_l.
\]
Thus the logarithms of the node tensor norms and external shadows
are exactly $P_*F(s)$. Changes of frames translate these formulas
compatibly and do not change their intrinsic meaning.

\begin{theorem}[Proved: the broader toric domain, its exact image, and its extra fibers]
\label{A1-toric-domain}
The tensor and external-shadow map from $\mathcal R_{\Gamma,L}(Y)$
has exactly the same image $\mathcal V_{\Gamma,L}(Y)$ as the
equal-length geodesic map of
Theorem~\ref{A1-bordered-shadow-range}; for $L=\varnothing$ this is
$\mathcal U_\Gamma(Y)$ of Theorem~\ref{A1-graph-range}. In particular
\[
 P_*(D_Y)=\left\{y:\ p\cdot y<K_{\Gamma,L}(p)
                   \text{ for all }0\ne p\in[0,\infty)^P\right\}.
\]
This is a proper subset of the unrestricted tensor/shadow fiber
whenever $P\ne\varnothing$; when $P=\varnothing$ it is the
entire one-point fiber.
Before finite automorphism quotients, the fiber over a fixed tensor
and shadow tuple is an open convex set of dimension $|E|$ in its
full cap logarithms. More strongly, after orienting each edge, the
quantities
\[
 d_e=s_{h_e^+}-s_{h_e^-}\in\R
\]
give a global real-analytic identification
\[
 \mathcal R_{\Gamma,L}(Y)
       \simeq\mathcal V_{\Gamma,L}(Y)\times\R^E,
\]
under which the surgery map is projection and the geodesic submodel
is $d=0$. Thus its tangent map is surjective with kernel dimension
$|E|$; it is an isomorphism precisely when $E=\varnothing$.
If the normalization varies and the graph has total genus $g$ and
$n$ external markings, the broader toric source has real dimension
$6g-6+2n+|L|+|E|$, whereas the tensor/shadow target has dimension
$6g-6+2n+|L|$. Thus the dimension count in 67/077 applies to the
geodesic model, not to actual toric cuts of independent lengths.

The identification is intrinsic if $\R^E$ is regarded as the real
space of antisymmetric functions on the two branches of each edge.
Branch interchange changes the corresponding $d_e$ to $-d_e$,
and graph automorphisms also permute edges. The construction varies
real-analytically with the normalization and is equivariant.
Consequently it defines a real vector orbibundle over the geodesic
image when finite automorphisms are retained. Scaling $d$ to zero
is a fiber-preserving equivariant strong deformation retraction.
The broader toric surgery map therefore induces the same
fundamental-group isomorphism to the unrestricted tensor/shadow
target, although its tangent map has a nonzero kernel whenever
there is a cut.
\end{theorem}
\begin{proof}
For a point $x\in D_Y$ and $p\geq0$, the vector $P_*^Tp$ is
nonnegative, and its entries at the two branches of an edge both
equal $p_e$. The component support inequalities give
\[
 p\cdot P_*x=(P_*^Tp)\cdot x<K_{\Gamma,L}(p)
 \qquad(p\ne0).
\]
Conversely, the already proved equal-length gradient inverse in
Theorem~\ref{A1-bordered-shadow-range} lifts every point satisfying
these inequalities to $F(P_*^Tt)\in D_Y$. This proves equality of
the two images without an additional separation or closure
assumption. All tensor phases occur independently by
Lemma~\ref{A1-tensor}. If $P$ is nonempty, a coordinate unit
vector $p$ gives a finite upper bound on the corresponding
logarithmic coordinate, proving properness. The empty case is
the identity on its one-point fiber.

Fix the target tuple, hence its logarithmic vector $y$ and its
tensor phases. Its cap-logarithm fiber is
\[
 D_Y\cap P_*^{-1}(y).
\]
It is nonempty, convex, and open in the affine space
$P_*^{-1}(y)$, whose dimension is
$|H|-|P|=|E|$. The tensor and the two shadow circles uniquely
determine the toric boundary identification, so there are no
additional phase variables in this fiber.

We prove the stronger global parametrization rather than infer it
merely from contractibility. Prescribe any $d\in\R^E$, put
$a(d)_{h_e^\pm}=\pm d_e/2$ and $a(d)_l=0$, and consider
\[
 \Phi_d(t)=S(P_*^Tt+a(d)),\qquad
 \Omega_d=\{t_e>|d_e|/2,\ t_l>0\}.
\]
Its Hessian is $P_*\,d^2S\,P_*^T>0$, because $P_*^T$ is
injective. Its recession function in every nonnegative direction
is $K_{\Gamma,L}$: apply Lemma~\ref{A2b-recession} to the positive
starting vector $P_*^Tt_0+a(d)$ and direction $P_*^Tp$.

For every $y$ in the displayed strict support region,
$\Phi_d(t)-y\cdot t$ has bounded sublevels, by the same
escaping-sequence segment argument as in
Theorem~\ref{A2b-image}. The finite lower bounds defining
$\Omega_d$ ensure that every limiting recession direction is
nonnegative. Its lower semicontinuous closure attains a minimum.
A finite boundary point has at least one branch or external
squared length tending to zero. Along an inward diagonal the
derivative is the sum of the cap logarithms, each with positive
multiplicity; at least one tends uniformly to $-\infty$ by the
collar bound, while every remaining term has a finite upper bound.
Thus no boundary minimum is possible. The minimum is interior
and unique, and satisfies
\[
             \nabla\Phi_d(t)=P_*F(P_*^Tt+a(d))=y.
\]
This constructs exactly one cap system for every $(y,d)$. The
implicit function theorem applied to its positive Hessian gives
joint real-analytic dependence on $y,d$ and the normalization.
There is no nonsmoothness at $d_e=0$: near any solution the
arguments of $S$ remain strictly positive, and its defining
equation is analytic in $(t,d)$; the expression $|d_e|/2$
describes only the open domain of existence.

The unique seam phases prescribed by the tensors complete the
claimed diffeomorphism. If every $d_e=0$, equal circumferences
make the toric seams isometries and their hyperbolic collars glue
smoothly, as in Proposition~\ref{A1-retained-discs}. Conversely,
if every seam is a geodesic for the canonical metric of the welded
surface, every cut component inherits its unique canonical
hyperbolic metric, so the two circumferences agree at every edge.
Thus $d=0$ is exactly the geodesic submodel. This converse uses
all seams simultaneously; a single geodesic among other
nongeodesic cuts does not justify the restriction argument.

Lengths are intrinsic, so any permitted normalization isomorphism
acts on $d$ by the stated signed edge permutation. Uniqueness of
the inverse makes the whole construction equivariant, including
self-loops. The homotopy $(q,r,d)\mapsto(q,r,(1-u)d)$ consequently
descends to the quotient groupoids and their coarse spaces. Its
composition with the common radial retract of
Theorem~\ref{A1-bordered-shadow-range} proves the
fundamental-group assertion. The vector-orbibundle description
does not identify a coarse quotient fiber with a smooth vector
space when a nontrivial finite stabilizer acts on it.
Finally, the marked normalization stratum has real dimension
$6g-6+2n-2|E|$. Adding $2|E|+|L|$ independent cap lengths and
$|E|$ seam phases gives the asserted source dimension; the
geodesic section removes precisely $|E|$ length differences.
\end{proof}

\begin{proposition}[Refuted: toric compatibility alone implies an inverse]
\label{A1-toric-noninjective}
Even for the unbordered four-labelled sphere with one separating
analytic cut, the torically compatible surgery map is not
injective. There are distinct such cut objects with fixed
normalization and tensor $q=\partial_z|_0\otimes\partial_w|_0$.
\end{proposition}
\begin{proof}
Let $P(r)$, $0<r<4$, be the uniquely determined marked hyperbolic
pant with two cusps and one geodesic boundary whose canonical cap,
on its normalized sphere $(0;1,\infty)$, has conformal radius $r$
at $0$. Existence, uniqueness, and strict variation of its
boundary length $\ell(r)$ follow from the one-active-cap case of
Theorem~\ref{A2b-image} and the Koebe endpoint in
Corollary~\ref{A2-explicit-limits}. Normalize its cap embedding
$f_r:\overline\D\to\mathbf P^1$ by $f_r(0)=0$ and
$f_r'(0)=r>0$. This removes the rotational ambiguity.

For each of the following two ordered pairs of pants, glue their
boundaries by the explicit toric rule
\[
 f_r(e^{i\theta})\sim f_{r'}(e^{-i\theta}):
 \qquad (r,r')=(1,1),\qquad (r,r')=(2,1/2).
\]
The cap maps extend holomorphically across their analytic boundaries,
so the rule gives analytic conformal welding. Label the two cusp
points on each side separately, so that any isomorphism of cut
objects preserves the two sides. Both outputs have normalization
equal to two spheres with their node at $0$ and the labelled
markings at $1,\infty$. By Lemma~\ref{A1-tensor}, their tensor is
in both cases $rr'\,\partial_z\otimes\partial_w
=\partial_z\otimes\partial_w$.

They are nevertheless distinct cut objects: a conformal
isomorphism preserving the cut restricts to a conformal isomorphism
of the first cut components, and therefore preserves their unique
canonical hyperbolic boundary lengths. Those lengths are
$\ell(1)$ and $\ell(2)$, which are unequal by strict monotonicity.
The first seam is geodesic. The second cannot be geodesic, since
its two canonical cut-component lengths $\ell(2)$ and $\ell(1/2)$
are unequal. This argument distinguishes the actual cut objects;
it does not assume that forgetting their cuts gives distinct
underlying four-pointed surfaces.
\end{proof}

\begin{corollary}[Proved: what changes when full external discs are retained]
\label{A1-toric-full-discs}
For the broader toric interpretation, retain arbitrary pointed
external discs with disjoint analytic closures, in addition to
the tensors. Such a tuple $(Y,q,(D_l))$ is a surgery output if and
only if its external shadows $r_l=r(D_l)$ and tensors belong to
$\mathcal V_{\Gamma,L}(Y)$ and there exists $d\in\R^E$ for which
the unique canonical cap system constructed in
Theorem~\ref{A1-toric-domain} has external members exactly $D_l$.
Its fiber is precisely this set of $d$'s, with the prescribed
finite automorphism identifications. In particular the
$\R^E$-fiber assertion concerns the tensor/shadow target; it
must not be transferred to a target retaining full external
disc shapes. The geodesic criterion of
Corollary~\ref{A1-bordered-shape-image} is the special case $d=0$.
\end{corollary}
\begin{proof}
For fixed $(Y,q,r,d)$ the preceding theorem constructs all caps and
their toric identifications uniquely. Necessity follows by cutting
a source object and recording its individual squared lengths.
Conversely the stated equality of external members makes the
constructed source map to exactly the prescribed full-disc tuple.
No claim that the external subsystem remains admissible after
discarding the node caps is used; that distinct condition still
has the obstruction in Proposition~\ref{A1-bordered-landing}.
\end{proof}
